\begin{statementbox}
	\textbf{\textcolor{CStatement}{条件}}
	\begin{description}[style=nextline, leftmargin=2.8em, labelsep=0.5em, itemsep=0.35em]
		\item[\textbf{\textcolor{CStatement}{条件 1（区间换元法）：}}] 函数 $y = A\sin(\omega x + \varphi)$（$A>0,\omega>0$），自变量 $x \in [x_1,x_2]$。令 $t = \omega x + \varphi$，则 $t \in [t_1,t_2]$，其中 $t_1 = \omega x_1 + \varphi$，$t_2 = \omega x_2 + \varphi$。
	\end{description}

	\textbf{\textcolor{CStatement}{结论}}
	\begin{description}[style=nextline, leftmargin=2.8em, labelsep=0.5em, itemsep=0.45em]
		\item[\textbf{\textcolor{CStatement}{结论一（最值公式）：}}]
			\textbf{\textcolor{CStatement}{直观描述：}}
			通过换元将原函数最值转化为正弦函数在区间上的最值。
			\[
				y_{\max}=A\cdot M,\quad y_{\min}=A\cdot m,
			\]
			其中 $M=\max_{t\in[t_1,t_2]}\sin t$，$m=\min_{t\in[t_1,t_2]}\sin t$。

		\item[\textbf{\textcolor{CStatement}{推广结论（余弦情形）：}}]
			\textbf{\textcolor{CStatement}{直观描述：}}
			余弦函数与正弦函数最值求法完全类似，只需将正弦换为余弦。
			\[
				y_{\max}=A\cdot M',\quad y_{\min}=A\cdot m',
			\]
			其中 $M'=\max_{t\in[t_1,t_2]}\cos t$，$m'=\min_{t\in[t_1,t_2]}\cos t$。
	\end{description}

	\textbf{\textcolor{CStatement}{等号/取等条件：}}
	对于正弦型，$y_{\max}$ 取等 $\Leftrightarrow \sin t = M$，$y_{\min}$ 取等 $\Leftrightarrow \sin t = m$；对于余弦型类似。
\end{statementbox}
